\documentclass{article}
\usepackage{tcee-repaso}
\tcetema[Los economistas clásicos y Marx]{3.A.2}{Los economistas clásicos y Marx.}
\tcefecha{04/11/2026}
\begin{document}
\cabecerarepaso
\begin{columnas}

% ==========================================================================
\bloque{Esquema del cante}
\esquemacante{%
\apartadocante{Introducción: «cada uno habla de la feria\ldots»}{2,5}
\apartadocante{I. Caracterización general de la escuela}{5,5}
\subcante{1. Según \autor{Schumpeter}}{1,5}
\subcante{2. Autores y rango histórico}{1,5}
\subcante{3. Ideas: rasgos, 4 leyes, dinero}{1,5}
\subcante{4. Método}{1}
\apartadocante{II. Los economistas clásicos}{13}
\subcante{1. Adam \autor{Smith} (1723-1790)}{3,5}
\subcante{2. Jean-Baptiste \autor{Say} (1767-1832)}{1}
\subcante{3. David \autor{Ricardo} (1772-1823)}{4,5}
\subcante{4. Thomas R. \autor{Malthus} (1766-1834)}{2}
\subcante{5. John Stuart \autor{Mill} (1806-1873)}{2}
\apartadocante{III. Karl \autor{Marx} (1818-1883)}{8}
\subcante{1. Contexto, método y obra}{1,5}
\subcante{2. Aportaciones}{5}
\subcante{3. Valoración y legado}{1,5}
\apartadocante{Conclusión}{1}}
\emph{Reparto}: \autor{Marx} $\geq$8 min; \autor{Ricardo} (modelo) $\geq$2 min. \emph{Pizarra}: \emph{izda.}, refrán y pegatinas («aprendiz de mucho\ldots», «quien tiene tierra\ldots», «donde comen dos\ldots», «trabajar para el obispo», «cría cuervos\ldots»); \emph{centro}, el gráfico de cada bloque; \emph{dcha.}, hasta el final, reparto de \autor{Ricardo} y $W=c+v+\mathit{pl}$.

% ==========================================================================
\bloque[2,5]{Introducción}
\begin{itemize}
\item \textbf{Enganche}: «Cada uno habla de la feria según le va en ella». Tréveris, 5/5/2018: estatua de \autor{Marx} de 5,5 m, regalo del Gobierno chino; Glasgow, 2023: tricentenario de \autor{Smith}. Misma economía (Revolución industrial), visiones opuestas: interés propio $\Rightarrow$ riqueza para todos / el capitalista se queda con parte de lo que produce el obrero. Pregunta: \emph{¿de dónde sale la riqueza de una nación y quién se la queda?}
\item \textbf{Relevancia}: \emph{teoría}: crecimiento \vertema{3.A.43} \vertema{3.A.44}, reparto \vertema{3.A.45}, comercio \vertema{3.B.5}, dinero \vertema{3.A.35}; \emph{política}: libre comercio (Leyes de Granos derogadas en 1846), patrón oro y Ley bancaria de 1844 \vertema{3.B.20}, límites del crecimiento (\autor{Malthus} $\to$ Club de Roma \vertema{3.B.31}); \emph{temario}: lo que corrigen los neoclásicos \vertema{3.A.3}, discute \autor{Keynes} \vertema{3.A.4} y critica \autor{Marx}; y 2 motivos de \autor{Schumpeter} (1954): no quedarse con el último manual; por qué cada idea llega en su época.
\item \textbf{Contextualización}: objeto \vertema{3.A.1}; ciencia desde 1776. Antes, pensamiento pero no ciencia: escritos no solo económicos, normativos y parciales. Preparan el terreno: \emph{mercantilistas} (riqueza = metales; protección), \emph{Escuela de Salamanca} (teoría cuantitativa: \autor{Azpilcueta}, 1556, antes que \autor{Bodin}, 1568), \emph{fisiócratas} (\autor{Quesnay}, \emph{Tableau}, 1758: circuito, \emph{laissez faire}). \autor{Smith} (1776) $\to$ \autor{Mill} (1848) $\to$ revolución marginalista (años 70). Rev. industrial; población que crece, tierra fija. «Clásicos»: \autor{Marx} (desde \autor{Petty}); \autor{Keynes} (1936) lo extiende a los neoclásicos.
\item \textbf{Problemática}: ¿qué comparten (objeto, método, ideas)? ¿De dónde sale el valor y cómo se reparte entre salarios, beneficios y rentas? ¿Crecer siempre o estado estacionario? ¿Puede faltar demanda? ¿Qué toma \autor{Marx} y por qué llega a lo opuesto?
\item \textbf{Estructura}: de la escuela a sus autores y de estos a su crítico: \hyperref[rep:I]{I}, caracterización; \hyperref[rep:II]{II}, \autor{Smith}, \autor{Say}, \autor{Ricardo}, \autor{Malthus}, \autor{Mill}; \hyperref[rep:III]{III}, \autor{Marx}.
\end{itemize}

% ==========================================================================
\bloque[5,5]{I. Caracterización general de la escuela\label{rep:I}}
\begin{ideaclave}
Escuela casi enteramente británica con un líder indiscutido (\autor{Smith}); objeto: crecimiento de la riqueza y su reparto entre clases; 4 leyes; el dinero solo fija los precios; método deductivo y verbal. \autor{Malthus} discrepa en la demanda; \autor{Marx} usa su aparato para criticarla.
\end{ideaclave}

\sub{I.1. Según \autor{Schumpeter}\label{rep:I.1}}
\begin{itemize}
\item Escuela = \emph{líder} + \emph{seguidores} + \emph{cuerpo doctrinal} (\autor{Schumpeter}, 1954). A diferencia de la neoclásica \vertema{3.A.3}, líder indiscutido.
\item \emph{Líder}: \autor{Smith}; sus preguntas son las de todos. \autor{Ricardo} corrige a \autor{Smith}; \autor{Mill} suaviza a \autor{Ricardo}.
\item \emph{Seguidores}: \autor{Say} (Francia); \autor{Ricardo}, James \autor{Mill}, \autor{Torrens}, \autor{Senior}, J.\,S. \autor{Mill}. \autor{Malthus}: clásico en población y renta, crítico en la demanda. \autor{Marx}: no clásico, su crítico (\hyperref[rep:III]{III}).
\item \emph{Objeto}: causas de la riqueza, su crecimiento a LP y su reparto entre 3 clases (trabajadores, capitalistas, terratenientes) en salarios, beneficios, rentas. Dinámico y agregado; sin micro/macro; el sujeto es la \emph{clase}. \autor{Ricardo}: «determinar las leyes que regulan esta distribución es el problema principal de la economía política».
\item \emph{Método}: \emph{economía política}, verbal, deductiva, con instituciones.
\end{itemize}

\sub{I.2. Autores y rango histórico\label{rep:I.2}}
\begin{itemize}
\item \emph{Periodo}: \emph{Riqueza} (1776) $\to$ \emph{Principios} de \autor{Mill} (1848), manual hasta \autor{Marshall} (1890). $\approx$70 años ($\approx$100 hasta 1871). \autor{Say}, gran excepción francesa.
\item \emph{3 grandes tratados}: \autor{Smith}, \emph{Wealth of Nations} (1776); \autor{Ricardo}, \emph{Principles of Political Economy, and Taxation} (1817); \autor{Mill}, \emph{Principles of Political Economy} (1848).
\item \emph{2 críticos}: \autor{Malthus} (se aparta en demanda y terratenientes); \autor{Marx} (valor trabajo, excedente y clases $\Rightarrow$ conclusiones opuestas).
\item \emph{Final} ($\approx$1870): caen población, valor trabajo y fondo de salarios (\autor{Mill} se retracta, 1869); \autor{Jevons} (1871) \vertema{3.A.3}.
\end{itemize}

\sub{I.3. Ideas económicas\label{rep:I.3}}
\begin{itemize}
\item \emph{3 rasgos} (\autor{Landreth}-\autor{Colander}):
  \begin{itemize}
  \item \emph{armonía}: con competencia y libertad, interés propio $\Rightarrow$ bien común (orden natural fisiócrata; mano invisible) \vertema{3.A.22}. Pero conflictos: terratenientes (\autor{Ricardo}); estado estacionario y subsistencia $\Rightarrow$ «ciencia lúgubre» (\autor{Carlyle}, 1849, a favor de la esclavitud; no por \autor{Malthus}). Armonía $\to$ ortodoxia; conflicto $\to$ heterodoxia (\autor{Marx});
  \item \emph{crecimiento a LP}: acumulación y reparto, no asignación ni ciclo; pleno empleo; precios y salarios flexibles; dinero solo fija precios \vertema{3.A.4}. Teoría \emph{dinámica} (neoclásicos: estática comparativa);
  \item \emph{desconfianza del gobernante} frente al mercantilista «médico»: \emph{laissez-faire} con excepciones (\autor{Mill}, más).
  \end{itemize}
\item \emph{4 leyes}:
  \begin{itemize}
  \item \concepto{rendimientos decrecientes de la tierra}: \concepto{margen intensivo} (más dosis en la misma tierra) y \concepto{extensivo} (tierras peores). \autor{Turgot} (1767); \autor{West}, \autor{Malthus}, \autor{Ricardo}, \autor{Torrens} a la vez en 1815 \vertema{3.A.11};
  \item \concepto{población} (\autor{Malthus}, 1798): crece si $w>$ subsistencia;
  \item \concepto{flujo-especie} (\autor{Hume}, 1752): superávit $\to$ entra oro $\to$ ↑ precios $\to$ fin del superávit \vertema{3.B.12};
  \item \concepto{ley de \autor{Say}}: no hay exceso de oferta \emph{general}.
  \end{itemize}
\item \concepto{Fondo de salarios}: CP, $w$ = capital adelantado / nº de trabajadores; LP, $w\to$ subsistencia. Juntas: acumulación $\to$ empleo $\to$ población $\to$ tierra peor $\to$ ↓ beneficio $\to$ estado estacionario.
\item \emph{Dinero}: no es riqueza, sino medio de circulación; teoría cuantitativa \vertema{3.A.35}.
  \begin{itemize}
  \item \emph{Controversia bullionista} (1797-1821): 1797 se suspende la convertibilidad; ↑ precio del oro y tipo de cambio. \emph{Bullionistas} (\autor{Thornton}, \autor{Horner}, \autor{Ricardo}): exceso de billetes; \emph{Bullion Report} (8/6/1810) les da la razón. \emph{Antibullionistas} (directores del Banco): la guerra; \concepto{letras reales}. Ley de 1819 (\autor{Peel}) $\to$ convertibilidad en 1821.
  \item \concepto{\emph{Currency school}} (\autor{Overstone}, \autor{Torrens}, \autor{Norman}): billetes uno a uno con el oro sobre una emisión fiduciaria fija; precursor \autor{Ricardo} (m. 1823, antes del debate). \concepto{\emph{Banking school}} (\autor{Tooke}, \autor{Fullarton}): la cantidad la decide el comercio; ley del reflujo; depósitos y letras también son dinero.
  \item Gana la 1.ª: \concepto{Ley de \autor{Peel}} (1844), departamentos de emisión y de banca; suspendida en 1847, 1857, 1866 \vertema{3.B.20}.
  \end{itemize}
\end{itemize}

\sub{I.4. Método\label{rep:I.4}}
\begin{itemize}
\item Economía política = \emph{ciencia moral} (\autor{Smith}, catedrático de filosofía moral); «política»: asunto del Estado, instituciones dentro. Ensayo en prosa, ejemplos numéricos, sin matemáticas.
\item \autor{Smith}: deducción + inducción e historia (Ilustración escocesa). \autor{Ricardo}: deducción pura; \concepto{vicio ricardiano} (\autor{Schumpeter}): aplicar a la política modelos muy simplificados. \autor{Malthus}: deducción y datos (2.ª ed., 1803). \autor{Senior} (1836): 4 proposiciones; \emph{ciencia} vs. \emph{arte}. \autor{Mill} (1836, recogido en 1844): ciencia abstracta del hombre que desea riqueza, \emph{a priori} $\to$ \emph{homo oeconomicus} (nombre posterior, de sus críticos). \autor{Marx}: materialismo histórico; leyes propias de cada modo de producción.
\item Positiva / normativa / arte: J.\,N. \autor{Keynes} (1891); en los clásicos, anacronismo útil.
\end{itemize}

% ==========================================================================
\bloque[13]{II. Los economistas clásicos\label{rep:II}}
\begin{ideaclave}
\autor{Smith}: riqueza = división del trabajo y mercado; \autor{Say}: ley de los mercados; \autor{Ricardo}: la tierra frena la acumulación y libre comercio; \autor{Malthus}: población y demanda; \autor{Mill}: producción con leyes, reparto según las instituciones.
\end{ideaclave}

\sub{II.1. Adam \autor{Smith} (1723-1790)\label{rep:II.1}}
\begin{itemize}
\item \emph{Vida}: Kirkcaldy (1723); catedrático de filosofía moral en Glasgow (1752-1764); amigo de \autor{Hume}; tutor del duque de Buccleuch (1764-1766: \autor{Quesnay}, \autor{Turgot}); comisario de aduanas (1778).
\item \emph{Obra}: \emph{Theory of Moral Sentiments} (1759): \emph{simpatía} y «espectador imparcial»; \emph{Wealth of Nations} (1776), 5 libros (división del trabajo, valor, reparto / capital / historia / mercantilismo y fisiócratas / finanzas del Estado). Líder por la síntesis, no por ideas nuevas.
\item \concepto{Mano invisible}: 3 veces (astronomía, «de Júpiter»; \emph{Sentimientos}, IV.1; \emph{Riqueza} \emph{una sola vez}, IV.2: «conducido por una mano invisible a promover un fin que no entraba en sus intenciones»). Metáfora secundaria (\autor{Rothschild}, 1994). 4 pasos: cada uno conoce su situación; «no es de la benevolencia del carnicero\ldots»; competencia $\Rightarrow$ precio de mercado $\to$ natural; se rompe con monopolio («la gente del mismo oficio rara vez se reúne\ldots»). Antecedente del 1TFEB \vertema{3.A.22}; fallos \vertema{3.A.23}.
\item \emph{Riqueza} = producto del trabajo, no oro. Depende de:
  \begin{itemize}
  \item \emph{productividad} ↑ con la \concepto{división del trabajo}: alfileres, 10 operarios, 18 operaciones, $>$48.000/día (4.800 por cabeza vs. $<$20). 3 causas: destreza, ahorro de tiempo, máquinas. «Aprendiz de mucho, maestro de nada», al revés;
  \item \emph{proporción de trabajo productivo} (se fija en una mercancía vendible) vs. \emph{improductivo} (criado, soberano, músico: útil, pero no deja producto).
  \end{itemize}
\item \emph{Límites}: \emph{extensión del mercado} (lib.~I, cap.~III) $\Rightarrow$ libre comercio; \emph{acumulación de capital} previa $\Rightarrow$ parsimonia. Antecedente del crecimiento endógeno (\autor{Young}, 1928) \vertema{3.A.44}.
\end{itemize}
\begin{itemize}
\item \emph{Valor}: \concepto{paradoja del agua y los diamantes} (lib.~I, cap.~IV; la resuelven los marginalistas \vertema{3.A.3}). \emph{Medida}: \concepto{trabajo comandado}. \emph{Causa}: estado primitivo $\to$ trabajo incorporado (castor = 2 ciervos); sociedad avanzada $\to$ \concepto{precio natural} = salarios + beneficios + rentas a tasas naturales; el \concepto{precio de mercado} (O y D efectivas) gravita hacia él \vertema{3.A.16}. Circular; renta componente (cap. VI) o efecto (cap. XI): \autor{Ricardo} toma lo 2.º.
\item \emph{Distribución}: salarios, CP fondo de salarios, LP natural, más alto si el país crece («recompensa liberal del trabajo»); beneficios incluyen interés + riesgo, ↓ con la acumulación; renta = precio de monopolio. \emph{Residuo: la renta} (en \autor{Ricardo}, el beneficio).
\item \emph{Ahorro}: se consume igual, por otros (trabajadores productivos): ley de \autor{Say} antes de \autor{Say}. \emph{Dinero}: «gran rueda de la circulación»; billetes convertibles; letras reales (antecedente de la \emph{banking school}).
\item \emph{Comercio}: no suma cero; \concepto{ventaja absoluta} \vertema{3.B.5}; amplía el mercado. Excepciones: defensa («más importante que la opulencia»: Acta de Navegación) y represalia. Límite: un país peor en todo no exportaría.
\item \emph{Política}: sistema de libertad natural, 3 deberes del soberano: defensa, justicia, obras e instituciones públicas no rentables (caminos, puertos, educación). 4 máximas fiscales: igualdad, certeza, comodidad, economía. Contra privilegios (compañías, gremios, aranceles).
\end{itemize}

\sub{II.2. Jean-Baptiste \autor{Say} (1767-1832)\label{rep:II.2}}
\begin{itemize}
\item Industrial, periodista, profesor. \emph{Traité} (1803): difunde a \autor{Smith}; valor = utilidad (precursor neoclásico); \emph{empresario} $\neq$ capitalista.
\item \concepto{Ley de los mercados} («Des débouchés», I.15): un producto abre salida a otros por todo su valor; productos se compran con productos; dinero, medio de cambio. Puede sobrar un bien, no todos. En inglés antes James \autor{Mill} (1808). «La oferta crea su propia demanda» es de \autor{Keynes} (1936) \vertema{3.A.4}.
\item \autor{Lange} (1942), sin el dinero \vertema{3.A.21}:
\begin{formula}
\centering $\displaystyle\sum_{k=1}^{n}p_k\,z_k(\vect p)\equiv 0$
\end{formula}
\item[] $\neq$ ley de \autor{Walras} (incluye el dinero). \autor{Becker}-\autor{Baumol} (1952): \concepto{identidad de \autor{Say}} (a cualquier precio, dinero sin papel) vs. \concepto{igualdad de \autor{Say}} (solo en equilibrio; excesos temporales: J.\,S. \autor{Mill}).
\item \emph{General glut} (1819-1824): \autor{Sismondi} y \autor{Malthus} vs. \autor{Say}, \autor{Ricardo}, James \autor{Mill}; ganan estos, ortodoxia hasta \autor{Keynes} \vertema{3.A.42}.
\end{itemize}

\sub{II.3. David \autor{Ricardo} (1772-1823)\label{rep:II.3}}
\begin{itemize}
\item \emph{Vida}: agente de bolsa, se retira desde 1814, diputado 1819-1823; entra por el bullionismo (\emph{High Price of Bullion}, 1810). \emph{Contexto}: \concepto{Leyes de Granos} (1815); \autor{Peel} las deroga en 1846 (hambruna irlandesa) \vertema{3.B.20}. \emph{Obra}: \emph{Essay on\ldots{} a Low Price of Corn} (1815); \emph{Principles} (1817, 1819, 1821: la 3.ª añade la maquinaria). Primero en razonar con modelos.
\item \emph{Valor}: \concepto{trabajo incorporado} (directo e indirecto), también en la sociedad avanzada. Salvedades: la utilidad es esencial pero no la medida; bienes no reproducibles (escasez); capital fijo de distinta duración $\Rightarrow$ teoría del valor trabajo «del 93~\%» (\autor{Stigler}, 1958). \emph{La renta no entra en el precio}: lo fija la tierra peor; «el grano no es caro porque se pague una renta, sino que se paga una renta porque el grano es caro». Medida invariable: no la halla ($\to$ \autor{Sraffa}, 1960, mercancía patrón).
\end{itemize}

\textbf{Modelo de distribución} («quien tiene tierra tiene guerra»; \emph{Essay}, 1815, contra las Leyes de Granos):
\begin{itemize}
\item \emph{Supuestos}: 1 sector (grano = producto, capital y salario); 3 clases (terratenientes gastan en lujo; trabajadores consumen; capitalistas adelantan $K=\bar wL$ y ahorran $B$); rendimientos decrecientes; salario natural $\bar w$ (\autor{Malthus}); beneficio residual.
\item $Y=F(L)$, $\mathrm{PMe}_L=F/L>\mathrm{PMg}_L=F'(L)$. \concepto{Renta diferencial} (\autor{West}, \autor{Malthus}, \autor{Torrens}, \autor{Ricardo}, 1815): la última dosis no paga renta:
\funcionclave{Y=\underbrace{(\mathrm{PMe}_L-\mathrm{PMg}_L)L}_{\text{rentas}}+\underbrace{\bar w L}_{\text{salarios}}+\underbrace{(\mathrm{PMg}_L-\bar w)L}_{B}}
\item Tasa de beneficio, con $K=\bar wL$:
\begin{formula}
\centering $\dfrac{B}{K}=\dfrac{\mathrm{PMg}_L-\bar w}{\bar w}$
\end{formula}
\item \emph{Dinámica}: $B/K>0\Rightarrow$ acumulan $\to$ ↑ fondo $\to w>\bar w\to$ ↑ población $\to w=\bar w$; ↑$L\Rightarrow$ ↓$\mathrm{PMg}_L$: ↑ renta, ↓ beneficios. \concepto{Estado estacionario} en $L^*$: $\mathrm{PMg}_L(L^*)=\bar w$, $B/K=0$; todo el excedente es renta.
\end{itemize}
\graficorepaso[0.85\linewidth]{ricardo-distribucion}{Reparto en $L_0$ y estado estacionario en $L^*$ (\autor{Kaldor}, 1956; \autor{Ricardo} no dibujó gráficos). Con $F=9L-L^2/2$, $\bar w=3$: en $L^*=6$, renta 18, salarios 18, beneficios 0}
\begin{itemize}
\item \emph{Implicaciones}: terratenientes contra las demás clases. Retrasan el estado estacionario: \emph{progreso técnico} agrario y \emph{libre comercio} del grano (↓ renta, ↑ beneficio).
\item \emph{Extensiones}: \autor{Kaldor} (1956), \autor{Pasinetti} (1960); maquinaria (cap.~31, 1821); acumulación: freno en \autor{Ricardo}, motor en \autor{Marx}. \emph{Evidencia}: el estado estacionario no llegó (técnica, importación tras 1846; ↑ salario real). \emph{Valoración}: renta como diferencia = germen del análisis marginal; factor fijo + sin progreso técnico $\to$ \autor{Solow} \vertema{3.A.43}. Falla: infravalora la técnica; con 1 bien evita el problema del valor.
\item \emph{Say y dinero}: acepta la ley («nadie produce si no es para consumir o vender\ldots»): puede faltar beneficio, no demanda. Gran bullionista; patrón en lingotes (1816); banco nacional que separe emisión y crédito (1824): precursor de la \emph{currency school}.
\item \concepto{Ventaja comparativa} (cap.~7) \vertema{3.B.5}: trabajo por unidad: paño Inglaterra 100, Portugal 90; vino 120 y 80. Portugal, ventaja absoluta en los 2; Inglaterra, comparativa en paño:
\begin{formula}
\centering $\dfrac{a_X^A}{a_Y^A}=\dfrac{100}{120}<\mathrm{RRI}=\dfrac{p_X}{p_Y}<\dfrac{a_X^B}{a_Y^B}=\dfrac{90}{80}$
\end{formula}
\item[] Con RRI $=1$, Inglaterra obtiene con 100 un vino que le costaría 120. Abierto: dónde queda la RRI ($\to$ \autor{Mill}). Antes: folleto anónimo (1814) y \autor{Torrens} (1815).
\item \emph{Política}: libre comercio (Leyes de Granos); impuestos recaen sobre beneficios, mejor sobre la renta; guerra con impuestos, no deuda: duda de que la gente razone que son equivalentes $\to$ \autor{Barro} (1974), \concepto{equivalencia ricardiana} (nombre de \autor{Buchanan}, 1976) \vertema{3.A.39}.
\end{itemize}

\sub{II.4. Thomas Robert \autor{Malthus} (1766-1834)\label{rep:II.4}}
\begin{itemize}
\item Clérigo anglicano; 1.\textsuperscript{er} profesor de economía política (Haileybury, 1805); para \autor{Keynes}, «el primero de los economistas de Cambridge». Contra \autor{Godwin} y \autor{Condorcet} (la miseria, culpa de las instituciones). \emph{Essay on the Principle of Population} (1798, anónimo; 6 ed. hasta 1826; la de 1803, empírica); \emph{Rent} (1815); \emph{Principles} (1820): demanda efectiva.
\item \emph{Principio de la población} («donde comen dos, comen tres», al revés). Precedentes: \autor{Hume}, \autor{Smith}, \autor{Wallace}, \autor{Cantillon}. 2 postulados: alimentos necesarios; pasión entre los sexos constante. Población, \emph{progresión geométrica} (se duplica cada 25 años: 1, 2, 4, 8\ldots); alimentos, \emph{aritmética} (1, 2, 3, 4\ldots): en 225 años, «como 512 a 10».
\end{itemize}
\graficorepaso[0.72\linewidth]{malthus-poblacion}{Progresiones de \autor{Malthus} (cada periodo, 25 años); la brecha la cierran los frenos}
\begin{itemize}
\item \concepto{Frenos}: \emph{positivos} (↑ mortalidad: hambre, epidemias, guerras); \emph{preventivos} (↓ natalidad: vicio; desde 1803, \concepto{contención moral}: retrasar el matrimonio).
\item \emph{Implicaciones}: salario real $\to$ subsistencia (pieza del modelo de \autor{Ricardo}); leyes de pobres agravan: abolirlas poco a poco. \emph{Extensiones}: \autor{Darwin} y \autor{Wallace}; neomaltusianos, Club de Roma \vertema{3.B.31}.
\item \emph{Evidencia}: describe la \emph{trampa maltusiana} preindustrial \vertema{3.A.43}; pero en Inglaterra, 1700-1850, población ×3,3 y producción agraria ×3,1: ambas geométricas; la producción agraria per cápita cae (125 $\to$ 93), la cubren importaciones; ↑ salario real en el XIX sin frenos. \emph{Valoración}: no previó el progreso técnico \vertema{3.A.44} ni la transición demográfica; \autor{Kremer} (1993), \autor{Simon} (1981): más población, más ideas.
\item \concepto{Demanda efectiva} (1820): si se ahorra demasiado, exceso \emph{general} de oferta; los trabajadores no compran todo $\Rightarrow$ \emph{consumidores improductivos} (terratenientes, Estado, criados). Pierde ante \autor{Ricardo} (sin teoría del atesoramiento ni del interés). Defiende la protección del grano. \autor{Keynes} (1936, cap.~23; 1933) lo reivindica; precursor del subconsumo \vertema{3.A.42}.
\end{itemize}

\sub{II.5. John Stuart \autor{Mill} (1806-1873)\label{rep:II.5}}
\begin{itemize}
\item Hijo de James \autor{Mill}; crisis de 1826 (se aparta de \autor{Bentham}); Harriet \autor{Taylor}; diputado. \emph{Essays on Some Unsettled Questions} (1844), \emph{Principles} (1848), \emph{On Liberty} (1859). Tras 1848: liberalismo reformista. Consolida y suaviza a \autor{Ricardo}; precursor del bienestar («fallo de mercado»: \autor{Bator}, 1958) \vertema{3.A.23}.
\item \emph{Valor}: «la teoría está completa». 3 casos: cantidad fija (demanda), coste constante (coste), coste creciente (coste en el margen).
\item \emph{Producción vs. distribución}: la 1.ª «participa del carácter de las verdades físicas»; la 2.ª «es solo cuestión de instituciones humanas» $\Rightarrow$ reformas.
\item \emph{Fondo de salarios}: lo mantiene (sindicatos no suben el salario de todos); lo abandona en 1869 (reseña de \autor{Thornton}).
\item \emph{Estado estacionario deseable}: no implica estancamiento del «progreso humano».
\item \emph{Say}: en trueque, sin exceso general; con dinero, transitorio (exceso de demanda de dinero): \emph{igualdad}, no identidad. \emph{Dinero}: flujo-especie + tipo de descuento \vertema{3.B.12}; \emph{banking school}: Ley de \autor{Peel} agrava las crisis.
\item \concepto{Demanda recíproca}: la RRI iguala exportaciones e importaciones; gana más quien menos necesita al otro; \autor{Marshall} (1879), curvas de oferta recíproca \vertema{3.B.5}. \concepto{Industria naciente}: aranceles temporales en una nación joven \vertema{3.B.8}.
\item Estado: funciones necesarias y facultativas (educación, menores). Impuestos proporcionales con mínimo exento; progresivos en herencias. Voto femenino, cooperativas.
\end{itemize}

% ==========================================================================
\bloque[8]{III. Karl \autor{Marx} (1818-1883)\label{rep:III}}
\begin{ideaclave}
El beneficio es trabajo no pagado (plusvalía); la competencia obliga a acumular y mecanizar $\Rightarrow$ ↓ tasa de ganancia y crisis cada vez más graves. Su teoría del valor no resiste el paso a los precios (transformación), pero sus preguntas siguen vivas.
\end{ideaclave}

\sub{III.1. Contexto, método y obra\label{rep:III.1}}
\begin{itemize}
\item Tréveris; doctor en filosofía (1841); París (1843), \autor{Engels} (1844); expulsiones; Londres desde 1849 (Museo Británico).
\item \emph{3 fuentes}: dialéctica de \autor{Hegel} + materialismo de \autor{Feuerbach}; socialismo francés; economía política inglesa (\autor{Ricardo}). \emph{Contexto}: obreros en ciudades, jornadas larguísimas, crisis (1825, 1836, 1847, 1857).
\item \emph{Obra}: \emph{Manifiesto} (con \autor{Engels}, 1848); \emph{Contribución} (1859, prólogo); \emph{El capital}: libro~1.º (1867); 2.º (1885) y 3.º (1894), por \autor{Engels}; \emph{Teorías de la plusvalía} (1905, \autor{Kautsky}).
\item \emph{Método}: de lo abstracto (mercancía) a lo concreto (precios, renta, crédito); leyes históricas, no naturales.
\end{itemize}

\sub{III.2. Aportaciones\label{rep:III.2}}
\begin{itemize}
\item \emph{Materialismo histórico}: «no es la conciencia\ldots{} sino su ser social el que determina su conciencia» (1859). \concepto{Fuerzas productivas} (medios de producción + fuerza de trabajo); \concepto{relaciones de producción} (propiedad, reparto) = \concepto{estructura económica}, base de la superestructura; \concepto{modo de producción}: asiático, antiguo, feudal, burgués moderno. Motor: choque fuerzas-relaciones $\to$ revolución social; lucha de clases (burguesía vs. proletariado). \concepto{Fetichismo de la mercancía}.
\end{itemize}
\begin{itemize}
\item \emph{Valor y plusvalía}. Pregunta: si se cambian equivalentes, ¿de dónde sale el beneficio? Mercancía: valor de uso y valor. Valor = \concepto{tiempo de trabajo socialmente necesario}. Trabajo \emph{concreto} (valores de uso) y \emph{abstracto} (valor). Libro I: cambio por su valor.
\item Se vende la \concepto{fuerza de trabajo}, no el trabajo; vale lo necesario para reproducirla (elemento histórico y moral), pero crea más: la \concepto{plusvalía}. Jornada de 10 h: 5 de \emph{trabajo necesario} y 5 de \emph{plustrabajo} («trabajar para el obispo»).
\item \concepto{Capital constante} $c$ (máquinas, materias: trabajo muerto, transfiere valor) y \concepto{variable} $v$ (salarios: el trabajo vivo repone su valor y añade $\mathit{pl}$):
\funcionclave{W=\underbrace{c}_{\text{trabajo muerto}}+\underbrace{v+\mathit{pl}}_{\text{trabajo vivo}}}
\item \concepto{Tasa de plusvalía}, \concepto{composición orgánica} y \concepto{tasa de ganancia}:
\begin{formula}
\centering $\mathit{pl}'=\dfrac{\mathit{pl}}{v},\quad \dfrac{c}{v},\quad g'=\dfrac{\mathit{pl}}{c+v}=\dfrac{\mathit{pl}'}{1+c/v}$
\end{formula}
\item[] Ej.: $\mathit{pl}'=5/5=100~\%$; con $c=80$, $v=20$: $g'=20~\%$. Con $c=0$, la $B/K$ de \autor{Ricardo}.
\item \emph{Plusvalía absoluta} (alargar la jornada) y \emph{relativa} (↓ trabajo necesario por productividad).
\end{itemize}
\graficorepaso[\linewidth]{jornada-laboral}{Jornada de 10 h: $\mathit{pl}'=100~\%$; absoluta (12 h): 140~\%; relativa (4 h necesarias): 150~\%}
\begin{itemize}
\item \concepto{Explotación}: no un robo en el intercambio, sino propiedad de los medios de producción.
\item \emph{Acumulación}: M--D--M (valor de uso) vs. D--M--D$'$ (plusvalía): sin límite; la competencia obliga \vertema{3.A.34}. \emph{Esquemas de reproducción} (libro II): sector I (medios de producción), II (consumo). Simple:
\begin{formula}
\centering $v_{\mathrm I}+\mathit{pl}_{\mathrm I}=c_{\mathrm{II}}$ \quad (ej.: $1.000+1.000=2.000$)
\end{formula}
\item[] Ampliada: $v_{\mathrm I}+\mathit{pl}_{\mathrm I}>c_{\mathrm{II}}$. Crecimiento equilibrado posible pero sin plan; \autor{Kalecki} (1933) $\to$ demanda efectiva antes que \autor{Keynes}.
\item \emph{Leyes del movimiento} («cría cuervos\ldots»: la burguesía produce sus sepultureros):
  \begin{itemize}
  \item \concepto{tendencia decreciente de la tasa de ganancia} (libro~3.º, caps.~13-15): mecanizar ↑$c/v\Rightarrow$ ↓$g'$. Contratendencias: más explotación, $c$ barato, salarios bajo su valor, ejército de reserva, comercio exterior; tope $g'<(v+\mathit{pl})/c$. \autor{Okishio} (1961): con técnicas que bajan costes y salario real dado, $g'$ no baja;
  \item concentración y centralización; \concepto{ejército industrial de reserva}; miseria creciente (relativa o absoluta); crisis cada vez más graves.
  \end{itemize}
\item \emph{Crisis}, 3 explicaciones sin unir: ↓ tasa de ganancia; \emph{realización} (subconsumo); \emph{desproporción}. Falla \autor{Say} \vertema{3.A.42}.
\item \emph{Política}: 10 medidas del \emph{Manifiesto} (expropiar la tierra, impuesto muy progresivo, abolir la herencia, crédito y transporte al Estado, educación gratuita); «de cada cual, según sus capacidades; a cada cual, según sus necesidades» (\emph{Gotha}, 1875).
\end{itemize}

\sub{III.3. Valoración y legado\label{rep:III.3}}
\begin{itemize}
\item \concepto{Problema de la transformación}: si solo el trabajo vivo crea plusvalía, más $c$ $\Rightarrow$ menos $g'$; pero la competencia iguala $g'$. \autor{Marx} (libro~3.º, cap.~9): \concepto{precio de producción} $=(c+v)(1+\bar g')$; 2 igualdades: $\sum$precios $=\sum$valores, $\sum$ganancias $=\sum$plusvalías.
\end{itemize}
\begin{center}\footnotesize\setlength{\tabcolsep}{3pt}
\begin{tabular}{@{}lcccccc@{}}
\toprule
\textbf{Ind.} & $c$ & $v$ & $\mathit{pl}$ & \textbf{Valor} & $g'$ & \textbf{Precio}\\
\midrule
I & 80 & 20 & 20 & 120 & 20~\% & \textbf{130}\\
II & 70 & 30 & 30 & 130 & 30~\% & 130\\
III & 60 & 40 & 40 & 140 & 40~\% & \textbf{130}\\
\midrule
Total & 210 & 90 & 90 & 390 & 30~\% & 390\\
\bottomrule
\end{tabular}\\[1pt]
{\color{tceegris}\scriptsize\itshape $\mathit{pl}'=100~\%$; precio $=(c+v)\cdot1{,}3$: por encima del valor con mucho $c$}
\end{center}
\begin{itemize}
\item \emph{Críticas}: \autor{Böhm-Bawerk} (1896): «gran contradicción» libro I / III; \autor{Bortkiewicz} (1907): dejó los insumos en valores; bien hecha, no se cumplen las 2 igualdades a la vez; \autor{Samuelson} (1971): valores, rodeo innecesario; \autor{Steedman} (1977, con \autor{Sraffa}): redundantes; plusvalía negativa con beneficios positivos (1975).
\item \emph{Lo que queda}: \autor{Morishima} (1973), \concepto{teorema marxiano fundamental}: ganancias $>0\Leftrightarrow$ explotación $>0$; \autor{Roemer} (1982, 1985): explotación por desigualdad en la propiedad.
\item \emph{Legado}: \autor{Kalecki}; \autor{Robinson} (1942) \vertema{3.A.4}; ciclo \vertema{3.A.42}. Acertó: concentración, sustitución de trabajo por capital, crisis recurrentes, reparto (\autor{Piketty}, 2014) \vertema{3.A.45}. Falló: ↑ salario real, clase media, no hubo derrumbe; planificación fracasó.
\end{itemize}

% ==========================================================================
\bloque[1]{Conclusión}
\begin{itemize}
\item \textbf{Recapitulación}: \autor{Smith}: riqueza del trabajo y su división, interés propio + competencia; \autor{Ricardo}: el reparto decide el crecimiento, tierra + \autor{Malthus} $\to$ estado estacionario; \autor{Mill}: no es catástrofe. \autor{Marx} le da la vuelta: beneficio = plusvalía; crisis internas.
\item \textbf{Relevancia}: reparto \vertema{3.A.45}, ventaja comparativa \vertema{3.B.5}, límites del crecimiento \vertema{3.B.31}.
\item \textbf{Extensiones}: marginalistas, agua y diamantes \vertema{3.A.3}; \autor{Keynes} vs. \autor{Say}, rescata a \autor{Malthus} \vertema{3.A.4}; nueva clásica \vertema{3.A.6}; \autor{Kalecki}.
\item \textbf{Opinión}: valor trabajo, talón de Aquiles; gran error, el pesimismo (infravaloraron la técnica); la pregunta de \autor{Ricardo} por el reparto vuelve con la desigualdad.
\item \textbf{Idea final}: la riqueza sale del trabajo, su división y la acumulación; quién se la queda, «según le va en la feria»: \autor{Smith} y \autor{Ricardo}, el capitalista; \autor{Malthus}, el terrateniente; \autor{Marx}, el obrero.
\end{itemize}

% ==========================================================================
\bloque{Autores y fechas}
\hito{1752}{\autor{Hume}: flujo-especie; 1758, \autor{Quesnay}: \emph{Tableau}}
\hito{1767}{\autor{Turgot}: rendimientos decrecientes}
\hito{1776}{\autor{Smith}: \emph{Wealth of Nations}}
\hito{1798}{\autor{Malthus}: \emph{Essay on Population} (2.ª, 1803)}
\hito{1803}{\autor{Say}: \emph{Traité}}
\hito{1810}{\autor{Ricardo}: \emph{High Price}; \emph{Bullion Report}}
\hito{1815}{Leyes de Granos; renta diferencial (\autor{West}, \autor{Malthus}, \autor{Torrens}, \autor{Ricardo})}
\hito{1817}{\autor{Ricardo}: \emph{Principles} (3.ª ed., 1821)}
\hito{1820}{\autor{Malthus}: \emph{Principles}}
\hito{1844}{Ley de \autor{Peel}; \autor{Mill}: \emph{Essays}}
\hito{1846}{Derogación de las Leyes de Granos}
\hito{1848}{\autor{Mill}: \emph{Principles}; \emph{Manifiesto}}
\hito{1863}{\autor{Lassalle}: «ley de hierro de los salarios»}
\hito{1867}{\autor{Marx}: \emph{El capital}, libro~1.º (2.º, 1885; 3.º, 1894)}
\hito{1871}{\autor{Jevons}: revolución marginalista}
\hito{1896}{\autor{Böhm-Bawerk}: «gran contradicción»}
\hito{1942}{\autor{Lange}: ley de \autor{Say}; \autor{Robinson}}
\hito{1956}{\autor{Kaldor}: modelo gráfico de \autor{Ricardo}}
\hito{1960}{\autor{Sraffa}: mercancía patrón}
\hito{1973}{\autor{Morishima}: teorema marxiano fundamental}
\hito{1974}{\autor{Barro}: equivalencia ricardiana}

% ==========================================================================
\bloque{Datos clave}
\begin{itemize}
\item Inglaterra, 1700-1850: población \dato{×3,3}{BoE, Broadberry}; producción agraria \dato{×3,1}{BoE}; producción agraria interior per cápita \dato{125 $\to$ 93}{BoE, años 70 del XVIII a 1850}.
\end{itemize}

% ==========================================================================
\bloque{Preguntas del tribunal}
\begin{enumerate}
\item \textbf{¿Residuo en \autor{Smith} y en \autor{Ricardo}?} En \autor{Smith}, la renta; en \autor{Ricardo}, el beneficio.
\item \textbf{¿Por qué la renta no entra en el precio?} El precio lo fija el trabajo en la tierra peor, que no paga renta: hay renta porque el grano es caro.
\item \textbf{¿Identidad o igualdad de \autor{Say}?} Identidad: a cualquier precio, dinero sin papel; igualdad: solo en equilibrio, excesos temporales (\autor{Mill}).
\item \textbf{¿Qué relaciona la tasa de ganancia de \autor{Marx} con la de \autor{Ricardo}?} $g'=\mathit{pl}/(c+v)$; con $c=0$, la $B/K$ de \autor{Ricardo}.
\item \textbf{¿Por qué difieren valores y precios en \autor{Marx}?} Por la distinta composición orgánica, con igual tasa de ganancia (transformación).
\end{enumerate}

\begin{trampa}
\begin{itemize}
\item Precio \emph{natural} de \autor{Smith} = costes; \emph{de mercado} = O y D (no al revés).
\item \autor{Smith}: medida = trabajo comandado; incorporado, solo en el estado primitivo.
\item \autor{Ricardo} \emph{fue} bullionista; murió (1823) antes del debate \emph{currency}/\emph{banking}.
\item Mano invisible: \emph{una sola vez} en la \emph{Riqueza}. «Ley de hierro de los salarios»: \autor{Lassalle} (1863); «clásicos»: \autor{Marx}; «la oferta crea su propia demanda»: \autor{Keynes}; «ciencia lúgubre»: \autor{Carlyle}.
\item \autor{Keynes} no dedicó la \emph{Teoría general} a \autor{Malthus}. El valor trabajo no lo inventa \autor{Marx}.
\item Falso que \autor{Mill} quisiera una economía más normativa que positiva. El siglo XIX no confirma a \autor{Malthus}.
\end{itemize}
\end{trampa}

\end{columnas}
\end{document}
